\chapter{Introduction}\label{sec:introduction}

Relay communication uses intermediate nodes between a source and a destination to extend coverage, improve reliability, or enable communication when a direct link is unavailable. Classical relay strategies such as amplify-and-forward (\texttt{AF}) and decode-and-forward (\texttt{DF}) provide well-established reference protocols. In several relay-channel models, coded block-based \texttt{DF} has strong information-theoretic guaranties under asymptotically long codewords \cite{CoverGamal1979RelayChannel,ElGamalKim2011NetworkIT,Laneman2004CooperativeDiversity,RankovWittneben2007HalfDuplexRelay,Nosratinia2004CooperativeCommunication}. The canonical comparison in this thesis is deliberately different: it is uncoded and uses symbol-wise relay processing, with the distinction between symbol-wise and information-theoretic block-\texttt{DF} made explicit in Remark~\ref{rem:df-terminology}.

Recent advances in machine learning have motivated neural processing for physical-layer tasks, including detection, decoding, equalization, and end-to-end communication-system optimization. A learned relay can approximate a nonlinear input--output mapping without executing an explicit analytic impairment model at inference. This does not make the relay model-agnostic: its behavior is determined by the channel conditions or channel family represented during training.

The central question of this thesis is therefore not whether learning-based relays are universally superior to classical methods, but \emph{under what conditions they provide a useful complement to them}. The experiments are organized to locate a boundary rather than to declare a universal winner.

Two considerations determine that boundary. The first is \emph{identifiability}: whether the information required to instantiate the applicable model-based classical rule is supplied or can be estimated reliably. The second is \emph{arithmetic and resource feasibility}: whether that rule can be executed within the available computation and decision-delay budget. Classical processing remains the preferred reference when its model is identifiable and affordable; learned relaying becomes a targeted complement when reliable per-block identification is unavailable or when exact model-based processing exceeds the available resource budget.

The experimental argument is organized as a four-layer information ladder:

\begin{enumerate}
    \item \textbf{Layer 1: matched memoryless control.}
    The canonical experiment uses a known, memoryless channel for which the receiver has the information required by the matched classical processing. This layer establishes the baseline comparison and primarily addresses H1--H4.

    \item \textbf{Layer 2: channel memory with fixed channel knowledge.}
    A finite-memory channel is introduced. The learned relay represents the fixed impairment through offline training, while the model-aware classical sequence detector is supplied with the corresponding channel taps. This layer separates the failure of memoryless relay functions from the failure of classical sequence detection.

    \item \textbf{Layer 3: finite-pilot identification.}
    Exact channel knowledge at the classical receiver is replaced by an estimate obtained from a finite pilot budget. This layer measures how the comparison changes as the reliability and overhead of per-block channel identification are reduced.

    \item \textbf{Layer 4: no per-block channel identification.}
    The channel realization is redrawn for each block from the same parametric family represented during training, but the realization itself is not supplied to the learned relay at inference. The learned relay therefore performs fixed, amortized within-family inference rather than online adaptation or generalization to a structurally unseen channel family.
\end{enumerate}

The first layer provides the matched-model control. Layers~2--4 form the principal unknown and mismatched-channel contribution and address H5. The supplementary coded/resource study is separate from this information ladder: it introduces coding, block-\texttt{DF}, soft-information forwarding, modulation-and-coding adaptation, buffering, and latency in order to test how those system-level considerations modify the canonical comparison.

\section{Scope and Contributions}\label{sec:scope-and-contributions}

The core comparison is carried out on one canonical setup, fixed in advance and not changed within that comparison: a half-duplex two-hop SISO link with one relay and no direct source--destination path, i.i.d.\ Rayleigh fast fading on both hops, complex-baseband Gray-coded QPSK, additive white Gaussian noise, coherent receiver-side processing, uncoded transmission, and BER as the performance metric. The relay processing function is the principal varied component.

\begin{longtable}{@{}p{0.25\textwidth}p{0.71\textwidth}@{}}
\caption[Scope of the thesis]{Scope of the thesis and the role of the canonical and extension studies.}
\label{tbl:scope}\\
\toprule
\textbf{Axis} & \textbf{Scope} \\
\midrule
\endfirsthead
\toprule
\textbf{Axis} & \textbf{Scope} \\
\midrule
\endhead
\bottomrule
\endlastfoot

\textbf{Topology}
&
One half-duplex relay, no direct source--destination path, and SISO transmission on both hops. Multi-relay, full-duplex, additional-hop, direct-link, and MIMO configurations are outside the evaluated scope.
\\

\textbf{Channel}
&
The canonical comparison uses independent i.i.d.\ Rayleigh fast fading on both hops with complex AWGN and no channel memory. The extension study deliberately introduces finite-memory and nonlinear impairments, finite-pilot channel estimation, and per-block-redrawn channel realizations from a trained parametric family.
\\

\textbf{Relay function}
&
Classical relay processing includes \texttt{AF} and symbol-wise \texttt{DF}; the learned comparison includes the neural relay architectures evaluated in Chapter~\ref{sec:experiments}. Sequence detection is introduced only when channel memory is explicitly added in Chapter~\ref{sec:unknown-channels}.
\\

\textbf{Modulation}
&
The canonical comparison uses Gray-coded QPSK on the Rayleigh channel. BPSK is confined to explicitly scoped calibration, companion, and unknown-channel experiments. Higher-order modulation appears in the supplementary modulation-and-coding study rather than as part of the canonical relay-architecture comparison.
\\

\textbf{Performance metric}
&
The canonical comparison uses uncoded BER versus SNR. Coding, goodput, modulation-and-coding adaptation, buffering, and latency are evaluated separately in the supplementary coded/resource study.
\\

\end{longtable}

The canonical experiment is intentionally a matched-model control. Because the physical channel is memoryless, the current received observation contains the information required for the current symbol decision after the receiver-side coherent processing defined in the system model. A learned relay may use a finite observation window as an architectural input, but that window does not create physical channel memory or provide additional information about neighboring independently transmitted symbols in the canonical experiment.

The canonical study establishes the baseline behavior of the relay strategies under conditions favorable to classical processing. In particular, it tests whether learned nonlinear processing can improve on linear \texttt{AF} at low SNR, whether matched symbol-wise \texttt{DF} remains competitive at medium and high SNR, how performance changes with model capacity, and how architecture comparisons behave when parameter budgets are controlled. These questions primarily address H1--H4.

The principal extension in Chapter~\ref{sec:unknown-channels} changes the information available to the relay and receiver rather than merely changing neural architecture. Channel memory first invalidates the memoryless relay assumption while leaving a correctly informed sequence detector applicable. Finite-pilot experiments then reduce the quality of per-block channel identification, and the final layer removes per-block identification while retaining training and testing within the same parametric channel family. This progression isolates the identifiability boundary underlying H5.

The resulting conclusion is deliberately conditional. Memory does not make classical processing obsolete: it invalidates memoryless relay functions while correctly informed sequence detection remains the stronger reference where its model is available and its computation is affordable. Conversely, a learned relay can become useful when the information required to instantiate that classical rule cannot be obtained reliably for each block.

A second boundary appears when the model is identifiable but exact processing becomes expensive. Exact trellis-based sequence detection grows rapidly with channel memory and constellation order, whereas the evaluated windowed learned relay grows more gently for a fixed architecture. This does not imply that learned complexity is independent of channel memory: preserving performance on a longer-memory channel may require a wider observation window or a different output representation. The detailed arithmetic and latency measurements are therefore deferred to Section~\ref{sec:joint-latency-memory}, where model identifiability and implementation feasibility are treated as separate but complementary constraints.

The supplementary coded/resource study in Chapter~\ref{sec:experiments} addresses a different question. It introduces a convolutional code, genuine block-\texttt{DF}, soft-information forwarding, modulation-and-coding adaptation, and explicit buffering and latency constraints. These experiments are not an additional layer of the information ladder; they test how coding and system-level resource decisions modify the standing of the relay strategies established by the canonical control.

Taken together, the thesis has three experimental components: the canonical matched-model control, the principal unknown- and mismatched-channel study, and the supplementary coded/resource study. Their combined purpose is not to establish a universally superior relay architecture. Rather, they delineate the conditions under which learned relay processing complements classical methods: classical processing dominates when the required model is identifiable and affordable, while learned processing becomes useful when reliable per-block identification or exact model-based processing becomes impractical.

The detailed communication model, receiver processing, relay implementations, training procedure, and evaluation methodology are defined in the corresponding methods sections. In particular, the exact coherent equalization convention is defined there rather than re-derived in this introductory chapter, and the distinction between uncoded symbol-wise \texttt{DF} and coded block-\texttt{DF} is maintained throughout the thesis.